\chapter{Fondamenti quantistici e circuiti variazionali}
\label{cap:fondamenti-quantistici}

Questo capitolo introduce i soli concetti quantistici necessari agli
esperimenti. L'obiettivo è rendere esplicite la rappresentazione dei dati, le
trasformazioni, le misure e il costo della simulazione da cui dipendono le
ipotesi della tesi.

\section{Stati, registri e notazione di Dirac}
\label{sec:hilbert}

Un qubit puro è un vettore unitario in $\C^2$,
$\ket{\psi}=\alpha\ket{0}+\beta\ket{1}$, con
$|\alpha|^2+|\beta|^2=1$. Un registro di $\nq$ qubit appartiene allo spazio di
Hilbert $\C^{2^{\nq}}$ e ha base computazionale
$\{\ket{0},\ket{1}\}^{\otimes\nq}$. Registri indipendenti si compongono tramite
prodotto tensoriale. La crescita esponenziale della dimensione è alla base sia
dell'argomento di espressività sia del costo della simulazione classica.

\section{Porte a uno e due qubit}
\label{sec:porte}

Le porte quantistiche sono operatori unitari. Gli esperimenti impiegano le
rotazioni $R_X(\theta)$, $R_Y(\theta)$ e $R_Z(\theta)$, la porta di Hadamard
$H$ e la CNOT. Le rotazioni agiscono su un singolo qubit e sono periodiche negli
angoli; Hadamard crea sovrapposizione, mentre CNOT correla due qubit e può
generare entanglement quando il controllo è in sovrapposizione. Un circuito è
la composizione ordinata di queste unità.

\section{Codifica angolare dei dati}
\label{sec:angle-embedding}

La codifica angolare associa una componente classica a una rotazione di ciascun
qubit:
\begin{equation}
  \ket{\psi(x)}=\bigotimes_{i=1}^{\nq}R(x_i)\ket{0}.
  \label{eq:angle-embedding}
\end{equation}
La scelta dell'encoding determina quali frequenze della funzione dei dati sono
accessibili al modello; ripetere le porte di encoding può ampliare questo
spettro \parencite{schuld2021encoding}. Nel circuito qui studiato la mappa è
applicata una sola volta, è periodica e non è iniettiva su un dominio
illimitato. Gli input neurali sono perciò limitati con $\tanh$ e riscalati.
Poiché viene codificata una componente per qubit, la dimensione in ingresso
coincide con $\nq$ e rende necessaria la compressione classica descritta nel
\cref{cap:software}.

\section{Ansatz variazionale: \texorpdfstring{\code{StronglyEntanglingLayers}}{StronglyEntanglingLayers}}
\label{sec:sel}

L'ansatz \code{StronglyEntanglingLayers} usa $\nqlayers$ strati, ciascuno
composto da una rotazione generica a tre parametri per qubit e da un anello di
porte entangling. I pesi hanno forma
\begin{equation}
  W\in\R^{\nqlayers\times\nq\times3}.
  \label{eq:weight-shape}
\end{equation}
L'anello permette che l'output di un qubit dipenda dalle componenti codificate
sugli altri. Nell'ablazione 03 esso e l'unitarietà distinguono strutturalmente
il VQC dal controllo classico limitato.

\section{Misurazione e valori di aspettazione}
\label{sec:misurazione}

L'uscita classica è il valore di aspettazione di Pauli-$Z$ sul qubit $i$:
\begin{equation}
  \expval{Z_i}=\bra{\psi}Z_i\ket{\psi}\in[-1,1].
  \label{eq:expval-z}
\end{equation}
La limitatezza è centrale per H3 ed è riprodotta dalla $\tanh$ finale della
baseline \code{bounded\_mlp}. Il simulatore restituisce il valore esatto dello
stato; su hardware lo stesso valore sarebbe stimato da un numero finito di
\emph{shot} e avrebbe incertezza di campionamento.

\section{Simulazione classica e differenziazione}
\label{sec:simulazione}

\code{default.qubit} mantiene in memoria $2^{\nq}$ ampiezze complesse e applica
classicamente le matrici delle porte. Memoria e tempo crescono quindi
esponenzialmente con $\nq$, sebbene per 4--10 qubit restino trattabili. L'accesso
allo stato intermedio consente \code{diff\_method="backprop"}, cioè la
retropropagazione attraverso il simulatore. Questo percorso non è disponibile
su hardware reale, dove i gradienti richiedono valutazioni ripetute, per esempio
tramite \emph{parameter-shift}. I tempi misurati non predicono pertanto il costo
su QPU.

\section{Barren plateaus}
\label{sec:barren-plateaus}

Nei \emph{barren plateau} la varianza del gradiente può decadere
esponenzialmente con il numero di qubit, rendendo l'ottimizzazione impraticabile
\parencite{mcclean2018barren}. Con $\nq\in[4,10]$ e $\nqlayers=2$ non si è nel
regime asintotico tipico; inoltre la diagnostica dell'esperimento 03 non mostra
gradienti quantistici sistematicamente più piccoli. Il fenomeno resta un limite
potenziale di scala, non una spiegazione dimostrata dei risultati.

\section{Algoritmi quantistici variazionali: quadro generale}
\label{sec:vqa}

I VQC appartengono alla famiglia degli algoritmi variazionali
\parencite{cerezo2021variational}: un circuito parametrico definisce una
funzione differenziabile ottimizzata da un algoritmo classico. Nei quantum
kernel, invece, la codifica definisce una mappa in uno spazio di Hilbert e la
fedeltà tra due stati funge da similarità \parencite{schuld2019quantum}.
L'elevata dimensione dello spazio non garantisce da sola separabilità utile;
inoltre alcuni quantum kernel possono concentrare i propri valori al crescere
del sistema, riducendo l'informazione discriminativa disponibile
\parencite{thanasilp2024concentration}. H4 verifica empiricamente la specifica
mappa IQP adottata in questa tesi.
